% !TeX root = ../papers/nc-p3.tex
\section{分岐を含む射影の次数七上界と二直線列の二重接触}
\label{sec:branch-degree-seven}
\begin{theorem}\label{thm:branch-degree-seven}
曲面射影のない標準正則窓の各辺・各方向で、評価の四次小行列式スキームの一次元部分の
総次数は重複込みで7以下である。従って定理\ref{thm:periodic-branch-obstruction}の
中心数と分岐を通る周期点列の長さの上界は、それぞれ14と28にできる。
\end{theorem}
\begin{proof}
一般平面上で評価行列の転置の余核を $E$ とし、
$0\to4\mathcal O(-1)\to6\mathcal O\to E\to0$ を取る。
一般点で階数4なので左は単射である。平面は孤立した階数2中心を避け、有限個の階数3点以外では
$E$ は束である。局所自由解消から閉点で深さ1以上なので $E$ は無捩れで、
$H=E^{**}$ は階数2の束、$H/E$ は有限長である。
階数3点では単元な三列を消去すると一本の三成分列が残り、
$\mathcal E xt^1(E,\mathcal O)$ はその三成分で割った局所環である。
これは四次小行列式の商環で、一般平面上の長さの和 $d$ が求める次数である。
$0\to E\to H\to H/E\to0$ の双対と局所双対性により、その長さは $H/E$ の長さにも等しい。
従って $c_1(H)=4,c_2(H)=10-d$ である。

$d\ge7$ とすると平面のRiemann--Roch公式で $\chi(H(-3))=d-6>0$、
また $h^2(H(-3))=h^0(H(-4))\le h^0(H(-3))$ なので $H(-3)$ は切断を持つ。
$H(-b)$ が非零切断を持つ最大の $b$ は3以上で、
飽和した切断から $0\to\mathcal O(b)\to H\to I_Z(4-b)\to0$ を得る。$Z$ は有限である。
$H$ は有限個の点の外で六つの定数切断で生成されるので $b\le4$。
$b=4$ なら $I_Z$ に非零切断があり $Z$ は空で、$6\mathcal O\to H\to\mathcal O$ が
評価行列の非零定数関係を与える。定理\ref{thm:branching-centers-curve}の証明で
一般平面への制限後にも六関係が独立であることに反する。
よって $b=3$、$10-d=3+\operatorname{length}Z\ge3$ である。
等号なら平面上で $H=\mathcal O(3)\oplus\mathcal O(1)$ となる。
全列の動く曲線は同じ位置では重複しないので、四位置で次数の和が28以下となる。
先の周期列の証明はこの和と各列の動く位置数が2以上であることだけを用いる。
\end{proof}

\begin{theorem}\label{thm:two-line-double-tangency}
曲面射影と純関係のない標準正則窓で分岐があれば、各位置の順中心と逆中心はそれぞれ高々一つである。
二直線列の初回直線・終端直線は両方とも系\ref{cor:two-line-secant-roots}の三次曲線の接線であり、
いずれの接触因子も一点の二倍である。
\end{theorem}
\begin{proof}
二直線列の始点を位置0とする。初回直線上の順欠損の二根が異なると仮定する。
位置1の二順中心から始まる両列の $d_2$ は3である（命題\ref{prop:pure-free-branch-cubic}）。
二列とも第三曲線が動き、その次数は $6-f_2$ である。
位置3に順中心がなければその二曲線は位置0にある相異なる六次曲線となり、次数上界に反する。
位置3の順中心が存在すれば、その後続直線も位置0にある別の曲線である。
先の二曲線とは一般前点が定点か動点かで区別される。従って一方でも次数6なら
総次数は少なくとも $6+1+1=8$ となる。よっていずれの位置1の列も位置3の順中心に到達する。
そこから位置0の元の中心に進める。実際、順中心の次位置にある全順中心はその後続直線上にある
（同命題のHankelの二者択一）。これで元の中心から異なる二つの初回点を通って戻る二経路ができ、
繰返すと指数増大になり矛盾する。逆向きにも同じで、二根の場合は双方とも排除された。

定理\ref{thm:branch-pencil-resultant}により二直線列は少なくとも一つ存在する。
その位置0は一中心である（二中心なら $d_2=3$）。位置1の順中心は直線上の今の二重根一つで尽きる。
一般に次位置の順中心が二つなら初回直線上の欠損が2以上で $d_2=1$ となり、
二つの異なる根を持つ二直線列を作ってしまう。従って位置2の順中心は高々一つである。
位置3に二つあればその次の位置0には順中心がないという同命題に反する。
従って全位置の主張が得られ、逆向きも同じである。
\end{proof}
